\section{Centered moments and their exact normalization}
\label{sec:foundation}

We first establish the analytic meaning of the moment values used in
the reconstruction theorem. The finite coordinate formulas are
inherited from the tail reports; their normalization is reproduced
here so that the additional determinant argument can be checked
without reconstructing an earlier archive.

\subsection{Centered asymptotic subtraction}

Put $x_n=n+1/2$. For an integer $p\ge2$, the centered Hurwitz
expansion is
\begin{equation}\label{found:asymptotic}
 A_p(n)\sim\sum_{j\ge0}a_j(p)x_n^{1-p-2j},\qquad
 a_0(p)=\frac1{p-1},\quad
 a_j(p)=\frac{B_{2j}(1/2)(p)_{2j-1}}{(2j)!}\quad(j\ge1).
\end{equation}
Here $(p)_m=p(p+1)\cdots(p+m-1)$ is the rising factorial.
For each fixed truncation the usual Euler--Maclaurin remainder
justifies this expansion, locally uniformly in any parameters kept
away from their poles \cite{DLMF}. In particular, if $p,q$ are even,
\begin{equation}\label{found:product-expansion}
 A_p(n)A_q(n)\sim
 \sum_{j\ge0}c_j(p,q)x_n^{2-p-q-2j},\qquad
 c_j(p,q)=\sum_{h=0}^j a_h(p)a_{j-h}(q).
\end{equation}

\begin{proposition}[Continued and ordinary centered moments]
\label{found:regularity}
For even $p,q\ge2$, $D_{p,q}(s)$ initially converges absolutely for
$\Re s>3-p-q$ and has a meromorphic continuation to the plane.
Its possible poles are among
\[
                         s=3-p-q-2j,\qquad j\ge0.
\]
Every $s=-2r$, $r\ge0$, is regular. Moreover, with
$J=r+1-(p+q)/2$, the exact ordinary identity is
\begin{equation}\label{found:ordinary}
 D_{p,q}(-2r)=\sum_{n\ge0}\left\{
 x_n^{2r}A_p(n)A_q(n)
 -\sum_{j=0}^{J}c_j(p,q)x_n^{2r+2-p-q-2j}\right\}.
\end{equation}
The inner sum is empty when $J<0$. The complete summand is
$O(x_n^{-2})$ when subtraction is needed; when it is not needed the
original series is absolutely convergent.
\end{proposition}

\begin{proof}
Subtract $L$ terms of \eqref{found:product-expansion}. Summing them
back gives
\begin{align}\label{found:continuation}
 D_{p,q}(s)={}&\sum_{n\ge0}
 \left(A_p(n)A_q(n)-\sum_{j=0}^{L-1}c_jx_n^{2-p-q-2j}\right)x_n^{-s}
 \notag\\
 &+\sum_{j=0}^{L-1}c_j\zeta(s+p+q-2+2j,1/2).
\end{align}
The first term converges normally for
$\Re s>3-p-q-2L$. Increasing $L$ gives the continuation and the
listed possible poles. Their positions are odd integers, so none
equals $-2r$. At this value choose exactly the nonnegative powers
in the polynomial subtraction. Each added-back Hurwitz value is
zero because
\[
 \zeta(-2h,1/2)=-\frac{B_{2h+1}(1/2)}{2h+1}=0
 \qquad(h\ge0).
\]
This proves \eqref{found:ordinary}. The statement uses a convergent
sum of the complete braces; the separately divergent pieces are
never summed on their own.
\end{proof}

The last vanishing is special to the chosen center. It also cannot
be differentiated as though it were an identity in the spectral
variable: for example $\zeta'(0,1/2)=-\tfrac12\log2\ne0$.
The shifted and logarithmic formulas below preserve the corresponding
finite correction terms.

\subsection{Finite summation and the block transform}

Let $\CT_N$ denote the coefficient of $N^0(\log N)^0$ in a
specified large-$N$ expansion. Centered even power sums satisfy
\begin{equation}\label{found:power-sum}
 S_r(N)=\sum_{n=0}^{N-1}(n+1/2)^{2r}
 =\frac{B_{2r+1}(N+1/2)}{2r+1}
 =\sum_{j=0}^r T_{rj}N^{2j+1},
\end{equation}
where
\[
 T_{rj}=\frac{\binom{2r}{2j}}{2j+1}B_{2r-2j}(1/2).
\]
The polynomial has zero constant term. By
Proposition~\ref{found:regularity}, the moment is consequently the
constant term of its finite partial sum. The recurrence
\[
                         A_p(k-1)=A_p(k)+k^{-p}
\]
and finite summation by parts give
\begin{equation}\label{found:blocks}
 D_{p,q}(-2r)=\sum_{j=0}^r T_{rj}C_{p,q}(2j+1),
\end{equation}
where, for a positive odd $d$,
\begin{align}\label{found:block-CT}
 C_{p,q}(d)=\CT_N\bigg\{&N^dA_p(N)A_q(N)
       +\sum_{k=1}^N k^{d-q}A_p(k)\notag\\
       &+\sum_{k=1}^N k^{d-p}A_q(k)
       +\sum_{k=1}^N k^{d-p-q}\bigg\}.
\end{align}
This identity is obtained at a finite cutoff before any limiting
operation. It is also the convenient starting point for shifted and
twisted extensions.

For later reference we state the explicit inherited evaluation.
Set $W=p+q$, and define
\begin{align}
 U_{p,q}(d)&=
 \begin{cases}
 \zeta(p,q-d)+\tfrac12\zeta(W-d),&d<q,\\[3pt]
 \displaystyle\sum_{h=1}^{\min\{(d-q+1)/2,p/2-1\}}
 \frac{\binom{d-q+1}{2h}}{d-q+1}
 B_{d-q+1-2h}\zeta(p-2h),&d>q,
 \end{cases}\label{found:U}\\
 R_{p,q}(d)&=
 \begin{cases}
 0,&d<W-1,\\[3pt]
 \displaystyle-\frac{B_{d-W+1}}2
 \left\{\frac1p\binom{d-q}{p-1}
       +\frac1q\binom{d-p}{q-1}\right\},&d\ge W-1.
 \end{cases}\label{found:R}
\end{align}
Then
\begin{equation}\label{found:C-evaluation}
 C_{p,q}(d)=U_{p,q}(d)+U_{q,p}(d)+R_{p,q}(d).
\end{equation}
All double-zeta sums appearing in \eqref{found:U} are convergent.
There is no case $d=q$ because $d$ is odd and $q$ is even.

\paragraph{Why these finite formulas hold.}
For $d<q$, changing the order in the convergent sum yields
$\sum_{k\ge1}k^{d-q}A_p(k)=\zeta(p,q-d)$.
For $d>q$ set $m=d-q$ and
$F_m(N)=(B_{m+1}(N+1)-B_{m+1})/(m+1)$. The finite identity
\[
 \sum_{k=1}^N k^mA_p(k)
 =A_p(N)F_m(N)+
                  \sum_{\ell=1}^N\frac{F_m(\ell-1)}{\ell^p}
\]
reduces everything to Bernoulli polynomials and ordinary power sums.
The even Bernoulli terms give \eqref{found:U}. The half-leading Faulhaber terms leave precisely
one half of $\zeta(W-d)$ for each convergent ordered block. When
both ordered blocks require reduction, their two half-terms cancel
the last sum in \eqref{found:block-CT}. At $d=W-1$ this includes cancellation of
the two $-\gamma/2$ constants against $+\gamma$ from the harmonic
sum. The remaining boundary products give \eqref{found:R}.
This is the finite-sum proof in the incoming tail calculus
\cite{Canonical,Bernoulli}; no additional relation among numerical
periods is used.

\subsection{Inverse transform and a small check}

Write $C_j=C_{p,q}(2j+1)$ and $D_r=D_{p,q}(-2r)$.
The Bernoulli generating identity in \eqref{found:power-sum} yields
the inverse transform
\begin{equation}\label{found:inverse}
 C_j=\sum_{r=0}^j\binom{2j+1}{2r}2^{-2(j-r)}D_r.
\end{equation}
One way to verify it is to set
$G(t)=\sum D_rt^{2r}/(2r)!$ and
$C(t)=\sum C_jt^{2j+1}/(2j+1)!$. The forward transform says
$G(t)=C(t)/(2\sinh(t/2))$ as a formal power series. Multiplying
back proves \eqref{found:inverse}, with no analytic assumption on
the growth of the moments.

For $p=q=2$, the first two blocks are $C_0=3\zeta(3)$ and
$C_1=-1/2$. Hence
\[
 D_0=3\zeta(3),\qquad
 D_1=-\frac14\zeta(3)-\frac16,
\]
and Proposition~\ref{found:regularity} gives the ordinary identity
\[
 \sum_{n\ge0}\bigl((n+1/2)^2\psi'(n+1)^2-1\bigr)
 =-\frac14\zeta(3)-\frac16.
\]
These are inherited checks. The next section uses the same fixed
normalization to prove an additional infinite reconstruction family.
